The Board of Trustees of the American Mathematical Society, expressing its belief that a great deal of time would be saved for mathematicians if they could study a textbook of Russian precisely adapted to their needs, granted to the present author nine months leave of absence from his duties as Editor of Translations. To the Board, and to Gordon L. Walker, the Executive Director of the Society, who took the initiative in this matter with his customary energy and good will, the author is deeply grateful for the opportunity to write such a book.

For indispensable help and advice in the preparation of the book, which was written chiefly in Göttingen, Moscow and Belgrade, gratitude is due to many people, especially to Martin Kneser of the Mathematics Institute in Göttingen, S. M. Nikol'skii and L. D. Kudrjavcev of the Steklov Institute in Moscow, T. P. Andjelić of the Mathematics Institute in the Yugoslav Academy of Arts and Sciences, G. Kurepa and B. Terzić of the Mathematics and Slavistics Departments in the University of Belgrade, and Alexander Schenker of the Department of Slavic Languages and Literatures in Yale University. For expert assistance, both secretarial and linguistic, the author is indebted to his wife Katherine and his son William, for proficient typing of the Reading Selections to Tamara Burmeister, Secretary of the Slavistics Department in Belgrade, and Christine Lefian, editorial assistant in the American Mathematical Society.

\vspace{1.5em}
\noindent\begin{tabularx}{\textwidth}{@{}Xr@{}}
Providence, USA & S. H. Gould\\
June, 1972 &
\end{tabularx}

\vspace{1.4em}
{\small
\noindent Most helpful among the books of reference have been:

\begin{flushleft}
Bruzguhova, E. A.: \emph{Practical phonetics and intonation of the Russian language} (Russian), Moscow, 1963\\
Spagis, A. A.: \emph{Formation and use of the aspects of the Russian verb} (Russian), Moscow, 1961\\
Wolkonsky, C. A. and Poltoratzky, M. A.: \emph{Handbook of Russian Roots}, New York, 1961\\
Vasmer, M.: \emph{Russisches etymologisches Wörterbuch}, Heidelberg, 1953\\
Daum, K. and Schenk, W.: \emph{Die russischen Verben}, Leipzig, 1954\\
Milne-Thomson, L. M.: \emph{Russian-English Mathematical Dictionary}, Madison: University of Wisconsin Press, 1962\\
Lohwater, A. J., with the collaboration of S. H. Gould: \emph{Russian-English Dictionary of the Mathematical Sciences}, Providence: American Mathematical Society, 1961.
\end{flushleft}
}
